\section{Fundamentação}

\textbf{Redes neurais.} Um neurônio artificial calcula uma soma ponderada das entradas, mais um viés, e aplica a ela uma função de ativação. Com a sigmoide, que leva qualquer número a um valor entre 0 e 1, um único neurônio é a regressão logística. O perceptron de múltiplas camadas (MLP) empilha camadas de neurônios entre a entrada e a saída, as camadas ocultas; com a ativação ReLU, $\max(0, z)$~\cite{nair:10}, elas permitem fronteiras de decisão não lineares. Os pesos começam sorteados e são ajustados pela retropropagação do erro~\cite{rumelhart:86}: a cada lote de exemplos, calcula-se quanto cada peso contribuiu para a perda (aqui, a entropia cruzada) e move-se o peso no sentido de reduzi-la, com um passo dado pela taxa de aprendizado. O otimizador Adam~\cite{kingma:15} adapta esse passo para cada peso. Uma época é uma passada completa pelo treino. Como o passo é o mesmo para entradas de escalas muito diferentes, os atributos costumam ser padronizados antes~\cite{lecun:98}. Treinar por épocas demais leva a rede a se ajustar a detalhes do treino; a parada antecipada~\cite{prechelt:98} separa uma parte do treino como validação e interrompe o treino quando o desempenho nela para de melhorar. Em dados tabulares, modelos de árvore ainda costumam superar redes neurais, em parte porque as redes tendem a aprender funções suaves e as árvores representam bem relações irregulares, como degraus~\cite{grinsztajn:22}.

\textbf{Modelos de referência.} O CART~\cite{breiman:84} divide o espaço de atributos de forma recursiva, em cada nó escolhendo o limiar de um atributo que mais reduz a impureza (Gini ou entropia); a vantagem é a legibilidade das regras. Uma árvore de profundidade 1 é uma \emph{regra de um corte}: um único atributo comparado com um único limiar. A Random Forest~\cite{breiman:01} treina muitas árvores, cada uma numa amostra com reposição do treino (\emph{bootstrap}) e com sorteio dos atributos de cada nó, e a previsão é a média dos votos. A regressão logística, como visto acima, equivale a uma rede de um único neurônio.
